\worksheettitle
  {Домашняя работа: ответы}
  {\modulemark{M02-H} Версия учителя}

\begin{teacherbox}[Что проверяет работа]
  Ученик должен выполнять оба направления преобразования, сохранять разряды с
  нулевыми цифрами и объяснять роль нуля. Проверяйте не только ответы, но и
  причины ошибок в задании 4.
\end{teacherbox}

\section{Ответы}

\begin{enumerate}[label=\textbf{\arabic*.}]
  \item Пропуски:
    \begin{enumerate}
      \item \(2\,000\), \(5\);
      \item \(400\);
      \item \(3\,000\).
    \end{enumerate}

  \item
    \begin{align*}
      48\,206&=40\,000+8\,000+200+6,\\
      70\,090&=70\,000+90,\\
      305\,004&=300\,000+5\,000+4,\\
      800\,060&=800\,000+60.
    \end{align*}

  \item \(45\,302\); \(90\,607\); \(520\,048\); \(700\,003\).

  \item
    \begin{enumerate}
      \item Цифра \(4\) стоит в разряде сотен:
        \(60\,405=60\,000+400+5\).
      \item Потеряны разряды единиц тысяч и десятков:
        \(80\,000+900+6=80\,906\).
      \item Цифра \(2\) стоит в разряде единиц тысяч:
        \(402\,070=400\,000+2\,000+70\).
    \end{enumerate}

  \item В числе \(50\,304\) цифра \(5\) имеет значение \(50\,000\), а в
    числе \(5\,304\) \textemdash{} \(5\,000\). Удаление нуля сдвигает цифры в другие
    разряды и изменяет число.

  \item \(604\,080=600\,000+4\,000+80\).

  \item Пример ответа: «…место каждой цифры определяет её разряд и значение;
    отсутствие нуля сдвинет следующие цифры в другие разряды».
\end{enumerate}

\begin{resultbox}[Критерий]
  \begin{itemize}
    \item \textbf{Результат устойчив:} задания 2, 3 и 6 выполнены верно, причины
      ошибок объяснены через разряды.
    \item \textbf{Нужна M02-R:} при сборке теряются нули или при разложении
      цифре систематически приписывается значение соседнего разряда.
  \end{itemize}
\end{resultbox}
